% SEGMENT OLP-0365-S01
% Part: lambda-calculus
% Chapter: introduction
% Section: beta

\documentclass[../../../include/open-logic-section]{subfiles}

\begin{document}

\olfileid{lam}{int}{bet}
\olsection{$\beta$-குறைத்தல்}

$(\lambd[m][(\lambd[y][y]) m])$ ஐக் காணும்போது, அதற்கும்
$\lambd[m][m]$ க்கும் ஏதோ தொடர்பு உள்ளது, குறிப்பாக இரண்டாம் சொல் முதலாவதை
“எளிமைப்படுத்துவதால்” கிடைக்க வேண்டும், என்று கருதுவது இயல்பானது.
\emph{$\beta$-குறைத்தல்} என்ற கருத்து இந்த உள்ளுணர்வை முறைப்படி எடுத்துரைக்கிறது.

\begin{defn}[$\beta$-சுருக்கல், $\bredone$] \ollabel{defn:betacontr}
  \emph{$\beta$-சுருக்கல்} ($\bredone$) என்பது சொற்களின்மீதான, பின்வரும்
  நிபந்தனையை நிறைவுசெய்யும் மிகச்சிறிய இசைவான தொடர்பாகும்:
  \[
    (\lambd[x][N])Q \bredone \Subst{N}{Q}{x}
  \]
  $P \bredone Q$ எனில் $P$ ஆனது~$Q$ ஆக \emph{$\beta$-சுருக்கப்படுகிறது}
  என்கிறோம். $(\lambd[x][N])Q$ வடிவிலான சொல் ஒரு \emph{ரெடெக்ஸ்}
  எனப்படுகிறது.
\end{defn}

\begin{prob} \ollabel{prob:def}
  \olref[lam][syn][alp]{defn:aconvone} இல் கட்டுண்ட மாறியின் மாற்றத்திற்குச்
  செய்ததுபோல, $\beta$-சுருக்கலின் சமானமான தொகுத்தறிதல் வரையறைகளை முழுமையாக
  எழுதுக.
\end{prob}

\begin{defn}[$\beta$-குறைத்தல், $\bred$] \ollabel{defn:betared}
  \emph{$\beta$-குறைத்தல்} ($\bred$) என்பது $\bredone$ ஐக் கொண்டிருக்கும்,
  சொற்களின்மீதான மிகச்சிறிய தற்சுட்டு, கடப்புத் தொடர்பாகும். $P \bred Q$
  எனில் $P$ ஆனது~$Q$ ஆக $\beta$-குறைக்கப்படுகிறது என்கிறோம்.
\end{defn}

சூழல் தெளிவாக இருக்கும்போது $\bredone$ க்குப் பதிலாக $\redone$ என்றும்,
$\bred$ க்குப் பதிலாக $\red$ என்றும் எழுதுவோம்.

முறைசாரா விதத்தில் கூறினால், பூச்சியம் அல்லது பல $\beta$-சுருக்கல் படிகளால்
$M$ ஐ~$N$ ஆக மாற்ற முடியுமானால், அப்போதும் அப்போது மட்டுமே $M \bred N$.

\begin{defn}[$\beta$-இயல்பானது]
மேற்கொண்டு $\beta$-சுருக்க முடியாத சொல் \emph{$\beta$-இயல்பானது}
எனப்படுகிறது.
\end{defn}

$M \bred N$ ஆகவும் $N$ ஆனது $\beta$-இயல்பானதாகவும் இருந்தால், $N$ ஐ~$M$ இன்
\emph{இயல்புவடிவம்} என்கிறோம். ஒரு சொல்லின் இயல்புவடிவம் தனித்ததா என்று
கேட்கலாம்; பின்னர் காணப்போவதுபோல், விடை ஆம்.

சில எடுத்துக்காட்டுகளைக் கருதுவோம்.
\begin{enumerate}
\item நமக்குப் பின்வருவது உள்ளது:
\begin{align*}
(\lambd[x][xxy]) \lambd[z][z] & \redone (\lambd[z][z])(\lambd[z][z]) y \\
& \redone (\lambd[z][z]) y \\
& \redone y
\end{align*}
\item ஒரு சொல்லை “எளிமைப்படுத்துவது” அதை உண்மையில் இன்னும் சிக்கலாக்கலாம்:
\begin{align*}
(\lambd[x][xxy])(\lambd[x][xxy]) & \redone (\lambd[x][xxy])(\lambd[x][xxy])y \\
& \redone (\lambd[x][xxy])(\lambd[x][xxy])yy \\
& \redone \dots
\end{align*}
\item அது ஒரு சொல்லை மாற்றாமலும் விடலாம்:
\[
(\lambd[x][xx])(\lambd[x][xx]) \redone (\lambd[x][xx])(\lambd[x][xx])
\]
\item மேலும் சில சொற்களை ஒன்றுக்கு மேற்பட்ட வழிகளில் குறைக்கலாம்;
  எடுத்துக்காட்டாக,
\[
(\lambd[x][(\lambd[y][yx]) z]) v \redone (\lambd[y][yv]) z
\]
என்பது மிகவெளிப்புறச் செயல்படுத்தலைச் சுருக்குவதால் கிடைக்கிறது; மேலும்
\[
(\lambd[x][(\lambd[y][yx]) z]) v \redone (\lambd[x][zx]) v
\]
என்பது மிகவுட்புறச் செயல்படுத்தலைச் சுருக்குவதால் கிடைக்கிறது. எனினும்,
இந்நிலையில் இரு சொற்களும் மேலும் அதே சொல் $zv$ ஆகக் குறைகின்றன என்பதைக்
கவனிக்க.
\end{enumerate}

கடைசி எடுத்துக்காட்டின் இறுதி விளைவு தற்செயலானதல்ல; மாறாக, சர்ச்--ரோசர் பண்பு
எனப்படும் லாம்டா கலனத்தின் ஆழமான, முக்கியமான ஒரு பண்பை அது விளக்குகிறது.

\begin{digress}
  பொதுவாக ஒரு சொல்லை $\beta$-குறைக்க ஒன்றுக்கு மேற்பட்ட வழிகள் உள்ளன.
  எனவே பல குறைத்தல் உத்திகள் உருவாக்கப்பட்டுள்ளன; அவற்றுள் மிகவும் பொதுவானது
  \emph{இயல்பான உத்தி}. ஒரு ரெடெக்ஸின் இடம், சொல்லில் அது தொடங்கும் புள்ளி
  என்று வரையறுக்கப்படும்போது, இயல்பான உத்தி எப்போதும் \emph{மிகஇடப்புற}
  ரெடெக்ஸையே சுருக்குகிறது. ஏதோ ஓர் உத்தியால் ஒரு சொல்லை இயல்புவடிவத்திற்குக்
  குறைக்க முடியுமானால், அப்போதும் அப்போது மட்டுமே, இயல்பான உத்தியால் அதை
  இயல்புவடிவத்திற்குக் குறைக்க முடியும் என்ற பயனுள்ள பண்பு இந்த உத்திக்கு
  உள்ளது. வேறுவிதமாகக் குறிப்பிடப்படாவிட்டால், தொடர்வருவதில் இயல்பான உத்தியையே
  பயன்படுத்துவோம்.
  % SOURCE CORRECTION TA-LAM-024: Tamil silently resolves the frozen spelling error “stratuegy” as “strategy.”
\end{digress}

\begin{defn}[$\beta$-சமானம், $\equal$]
  \emph{$\beta$-சமானம்} ($\equal$) என்பது பின்வருமாறு தொகுத்தறிதல் முறையில்
  வரையறுக்கப்படும் தொடர்பாகும்:
  \begin{enumerate}
  \item $M \equal M$.
  \item $M \equal N$ எனில், $N \equal M$.
  \item $M \equal N$, $N \equal O$ எனில், $M \equal O$.
  \item $M \equal N$ எனில், $PM \equal PN$.
  \item $M \equal N$ எனில், $MQ \equal NQ$.
  \item $M \equal N$ எனில், $\lambd[x][M] \equal \lambd[x][N]$.
  \item $(\lambd[x][N])Q \equal \Subst{N}{Q}{x}$.
  \end{enumerate}
\end{defn}

முதல் மூன்று விதிகள் இத்தொடர்பை ஒரு சமானத் தொடர்பாக ஆக்குகின்றன; அடுத்த
மூன்று அதை இசைவானதாக ஆக்குகின்றன; கடைசி விதி அது $\beta$-சுருக்கலைக்
கொண்டிருப்பதை உறுதிசெய்கிறது.

முறைசாரா விதத்தில் கூறினால், $\beta$-சுருக்கல் அல்லது
$\beta$-சுருக்கலின் “எதிர்ச்செயல்” ஆகியவற்றின் பூச்சியம் அல்லது பல படிகளால்
இரண்டு சொற்களில் ஒன்றை மற்றொன்றாக மாற்ற முடிந்தால், அப்போதும் அப்போது மட்டுமே
அவை $\beta$-சமானமானவை. $N$ ஆனது~$M$ ஆக $\beta$-சுருக்கப்படும் அப்போதும்
அப்போது மட்டுமே $M$ ஆனது~$N$ ஆக எதிர்-$\beta$-சுருக்கப்படும் வகையில்,
$\beta$-சுருக்கலின் எதிர்ச்செயல் வரையறுக்கப்படுகிறது.

மேலுள்ள விதிகளைத் தவிர, இத்தொடர்பை மேலும் பல விதிகளால் விரிவாக்குவோம்;
$X$ என்பது விரிவாக்கும் விதியாக இருக்கும்போது, விரிவாக்கப்பட்ட சமானத் தொடர்பை
$\equal[X]$ என்று குறிக்கிறோம்.

\end{document}
